% GENERATED FILE. Edit canonical lessons, then run npm run generate:books.
% canonical-source-hash: fnv1a64:ceed6b208e6e60b0
% canonical-lessons: TA-C09-mannikkavum, TA-C09-sorry-register

\chapter{Sorry}
\label{ch:ta-sorry}

This chapter is generated from the canonical micro-lessons used by Language Ladder.
Each section stays independently resumable and preserves the authored prerequisite order.

\section[\ta{மன்னிக்கவும்}]{\ta{மன்னிக்கவும்} (maṉṉikkavum) — please forgive}

\label{lesson:TA-C09-mannikkavum}



\begin{quote}
\textbf{Warm-up.} {[}PAUSE 2s{]} You already know \textbf{\ta{தயவுசெய்து}}, \textquotedblleft{}please.\textquotedblright{} Tamil's \textquotedblleft{}sorry\textquotedblright{} gives you something its Dravidian cousins don't: a forgiveness-word that's Tamil through and through, not borrowed from Sanskrit.
\end{quote}

\begin{cousinweb}
\textbf{\ta{மன்னிக்கவும்}} is built from three pieces:

\begin{itemize}
\raggedright
  \item \textbf{\ta{மன்னி}} (\emph{maṉṉi}) = \textquotedblleft{}\textbf{to forgive, excuse, pardon}\textquotedblright{} — and unlike the words you'll meet in Kannada, Telugu, and Malayalam, this is Tamil's \textbf{own} native verb, not borrowed from Sanskrit \textbf{kṣamā}.
  \item \textbf{‑\ta{க்க}} (\emph{‑kka}) = an infinitive-forming ending.
  \item \textbf{‑\ta{உம்}} (\emph{‑um}) = a formal, written-register particle used to make a soft, general request or permission — the same particle you'd see on a sign asking visitors to \emph{\ta{வரவும்}} (\textquotedblleft{}please come\textquotedblright{} / \textquotedblleft{}you may come\textquotedblright{}).
\end{itemize}

So \textbf{\ta{மன்னிக்கவும்}} reads as something like \textquotedblleft{}\textbf{may {[}one{]} forgive}\textquotedblright{} — a polite, slightly formal way of asking for forgiveness, more at home in writing or a considered apology than tossed off in passing.
\end{cousinweb}

\begin{cousinweb}
Where Kannada, Telugu, and Malayalam all reach for the Sanskrit noun \textbf{kṣamā} (\textquotedblleft{}patience, forgiveness\textquotedblright{}) and turn it into a verb, \textbf{Tamil uses its own root}:

\begin{itemize}
\raggedright
  \item Tamil \textbf{\ta{மன்னிக்கவும்}} \emph{maṉṉi\textbf{kkavum}} — from Tamil \textbf{\ta{மன்னி}} \emph{maṉṉi}
  \item Kannada \textbf{\ka{ಕ್ಷಮಿಸಿ}} \emph{kṣami\textbf{si}} — from Sanskrit \textbf{kṣamā}
  \item Telugu \textbf{\te{క్షమించండి}} \emph{kṣamin̄\textbf{caṇḍi}} — from Sanskrit \textbf{kṣamā}
  \item Malayalam \textbf{\ml{ക്ഷമിക്കണം}} \emph{kṣami\textbf{kkaṇam}} — from Sanskrit \textbf{kṣamā}
\end{itemize}

Where \textquotedblleft{}please\textquotedblright{} showed the whole family sharing one Sanskrit-derived idea (\emph{daya}), \textquotedblleft{}sorry\textquotedblright{} shows the opposite: Tamil kept its own word while its cousins borrowed the same Sanskrit one.
\end{cousinweb}

\subsection*{Guided Practice}
{[}PAUSE 1s{]}

\begin{itemize}
\raggedright
  \item {[}YOU SAY: \textquotedblleft{}maṉṉi\textquotedblright{} — forgive — the Tamil root, not from Sanskrit{]}
  \item {[}YOU SAY: the whole word — \textquotedblleft{}maṉṉikkavum\textquotedblright{} = please forgive{]}
  \item {[}YOU SAY: contrast — Kannada, Telugu, Malayalam all use kṣamā instead{]}
\end{itemize}

\begin{tcolorbox}[breakable,colback=teal!4,colframe=teal!35!black,title={Wrap-up recall}]
{[}PAUSE 3s{]} What does \textbf{\ta{மன்னிக்கவும்}} mean, and what is it built from? (\textbf{\textquotedblleft{}Please forgive\textquotedblright{}} — \emph{maṉṉi} \textquotedblleft{}to forgive\textquotedblright{} + \emph{‑kkavum}, a formal request ending.) What's different about Tamil's forgiveness-word compared to Kannada, Telugu, and Malayalam? (\textbf{It's Tamil's own native root} — the other three borrow Sanskrit \emph{kṣamā} instead.) Next: place the word in its real register and read the doubled alveolar n.
\end{tcolorbox}
\section[\ta{மன்னிக்கவும்} / sorry]{\ta{மன்னிக்கவும்} — formal forgiveness and a doubled n}

\label{lesson:TA-C09-sorry-register}



\begin{quote}
\textbf{Warm-up.} {[}PAUSE 2s{]} The root is native Tamil. Its everyday conversational fit is less simple than a one-to-one dictionary entry suggests.
\end{quote}

\begin{culture}
\textbf{\ta{மன்னிக்கவும்}} leans formal: a \textquotedblleft{}please forgive\textquotedblright{} on a notice or in a considered apology. Spoken Tamil may use shorter regional forms or borrow English \textquotedblleft{}sorry.\textquotedblright{} Colloquial variation is real, so treat formal/written as a reliable center, not a claim that one casual form covers every region.
\end{culture}

\subsection*{Script you'll notice: Read \ta{ன்னி}}
In \textbf{\ta{ன்னி}}, alveolar \textbf{\ta{ன}} is first silenced by puḷḷi and then repeated before the vowel. The spelling holds the consonant across the syllable boundary:

\begin{quote}
\textbf{\ta{ன்} + \ta{னி} $\to$ \ta{ன்னி}} — \emph{ṉ-ṉi}
\end{quote}

This is the same gemination-by-repetition principle used elsewhere in Tamil; the sound is doubled without inventing a new ligature.

\subsection*{Guided Practice}
{[}PAUSE 1s{]}

\begin{itemize}
\raggedright
  \item {[}YOU SAY: \textquotedblleft{}maṉ-ṉi-kkavum,\textquotedblright{} holding both doubled consonants{]}
  \item {[}YOU CHOOSE: considered/formal apology $\to$ \textbf{\ta{மன்னிக்கவும்}}{]}
  \item {[}YOU QUALIFY: casual alternatives vary by region{]}
\end{itemize}

\begin{tcolorbox}[breakable,colback=teal!4,colframe=teal!35!black,title={Wrap-up recall}]
{[}PAUSE 3s{]} What is the reliable register center of \emph{maṉṉikkavum}? (\textbf{Formal/written or considered apology}.) Must every region use it casually? (\textbf{No}.) How does \textbf{\ta{ன்னி}} show a doubled \emph{ṉ}? (Puḷḷi stops the first \textbf{\ta{ன்}}, then full \textbf{\ta{னி}} follows.)
\end{tcolorbox}
